\justifying
\NSFCBodyText

\subsubsection{已具备的实验条件}

\paragraph{计算与软件条件。}申请人已确认可稳定使用 NVIDIA A40 GPU 8 张和 RTX 3090 Ti GPU 6 张，并具备 Linux、Python、深度学习框架、Git 与 LaTeX 等软件和工程工具链。A40 可组成同构训练环境，承担约 1B 规模开源模型的成对继续预训练；3090 Ti 用于质量评分、去重检索、人工审计辅助、模型推理和小规模 smoke test，避免异构显卡混入同一同步数据并行作业。现有工具链能够记录数据版本、代码提交、模型配置、随机种子、日志与统计结果，支撑双来源冻结协议。

\paragraph{人员与任务条件。}申请人具备研究构思、模型设计、实验指导、结果分析和论文写作经验，负责总体科学问题、干预协议、统计口径与成果质量控制。七名已确认参与的新疆大学研究生计划在项目周期内参与：孔哲负责低资源文本整理、清洗与质量审计，张建宏负责条件效用建模与等预算样本选择实验，于得武负责代码整合、版本管理、实验日志与统计分析；杨铄、孟星宇、桑乐、吕中原分别具备数据处理、域泛化/适配、语言模型训练和负迁移消融等任务基础。后四类能力作为数据处理与稳健性分析的辅助条件，不将本项目重新扩展为 OCR、多语种或多模态研究。

\paragraph{数据与合规条件。}本项目不依赖私有业务数据或人体受试样本。已实际下载 ChineseWebText2.0 公开测试分片，复核数据卡所声明的 Apache-2.0 许可，并记录 1,000 条原始数据的文件哈希、筛选规则和随机种子；同时已固定 Qwen2.5-0.5B 的模型 revision 与权重哈希，并完成静态 loss 复核。这些资产证明小样本下载、版本固定、扰动—恢复与日志链路可运行，不代表来源 A/B 的全量许可审计已完成。正式语料仍须逐来源检查条款、去重、隐私/敏感内容和训练—评价隔离；不能再分发的数据只保留索引、配置和统计信息，按许可决定派生标注与模型产物的开放范围。

\subsubsection{尚缺少的实验条件及拟解决的途径}

目前已形成一个公开测试分片、初版扰动/恢复/sham 脚本、两类静态评分结果和完整运行日志；尚缺用于确认性训练的正式来源 A/B 语料索引、双人语体与语义标注、六臂平衡表、继续预训练配置、能力评价脚本和消融日志。第一年度将按“完整许可/来源核验—人工试标—段落/连接层六臂生成与 token/packing 平衡审计—10M token smoke test”的顺序补齐；静态预实验不作为继续预训练或机制证据。

目前没有已落实的语言学/语料库外部合作者或专属人工标注平台。项目主方案采用可操作的结构完整、语义保持、事实一致和语体属性标注，由团队开展双人独立标注与分歧复核；如后续获得正式合作，可将专家复核作为增强条件，但不以未确认合作作为项目完成前提。标注一致性不足时先缩小样本、修订规则并重复试标，而非直接扩大人力。

学院公共算力可按流程申请，但固定账号、设备型号和额度尚未落实；A100 八卡资源也未落实，均不计入稳定条件。主方案只按已确认的 A40 与 3090 Ti 设计。若排队或存储资源不足，取消模型尺度和状态探索、减少 checkpoint 长期保留并统一下调单次 token，保留来源 A 的六个成对主种子、sham 和支持窄化对照。

来源 B 及较大模型只承担外部边界验证。若候选公开语料在许可、独立 token 或来源独立性上不满足要求，来源 B 可降级为不同来源层的严格留出环境，并明确外部效度减弱；若同家族较大模型所需显存和节点时超出首年实测预算，则直接取消尺度探索，不影响双路径核心检验。

\subsubsection{利用研究基地的计划与落实情况}

本项目当前未以国家实验室、全国重点实验室或部门重点实验室的专属设备、数据与合作承诺作为实施条件，也暂无已落实的相关研究基地使用协议。核心研究依托申请人已确认可稳定使用的计算资源和团队完成。如项目实施期间通过正式申请获得校级或更高层级公共平台资源，仅用于补充训练时段、存储和结果复核，并在使用前记录授权范围和资源证明；未落实前不写入固定配额或既有合作。
